\documentclass[11pt, a4paper, UTF8]{ctexart}

% ==================== 宏包 ====================
\usepackage{geometry}
\usepackage{fontspec}
\usepackage{enumitem}
\usepackage{fancyhdr}
\usepackage{graphicx}
\usepackage{tabularray}
\usepackage{fontawesome5}
\usepackage{setspace}
\usepackage{ulem}
\usepackage{xcolor}     % tabularray 的 bg= 颜色需要此包
\usepackage{lastpage}   % 自动获取总页数，替代手动 \label{ReportEnd}

% 页面边距
\geometry{top=42mm, bottom=30mm, left=25mm, right=25mm, headheight=60pt}

% 字体
\setCJKmainfont{SimSun}[AutoFakeBold=true]
\setCJKsansfont{SimHei}[AutoFakeBold=true]
\setmainfont{Times New Roman}
\setsansfont{Arial}

% ==================== 【变量配置区 — 每次出报告只改这里】====================

% —— 报告基本信息 ——
\newcommand{\ReportYear}{20xx}          % 年份
\newcommand{\ReportMonth}{xx}           % 月份
\newcommand{\ReportDay}{xx}             % 日
\newcommand{\ReportNo}{xxxx}            % 报告编号

% —— 封面三行 ——
\newcommand{\ProjectName}{%             % 工程名称（过长时在此手动换行，见封面说明）
    小米智能制造产业基地二期项目 （6号生产厂房等2项）-6号生产厂房
}
\newcommand{\ClientName}{小米景曦科技有限公司}
\newcommand{\InspType}{一般委托}

% —— 主表字段 ——
\newcommand{\InspLocation}{北京市亦庄经济开发区}
\newcommand{\InspItem}{防火涂层厚度检测}
\newcommand{\ConstructUnit}{中建二局集团有限公司}
\newcommand{\DesignUnit}{中国汽车工业工程有限公司}
\newcommand{\SupervUnit}{北京中联环建设监理有限责任公司}
\newcommand{\SampleBasis}{委托方指定数量}
\newcommand{\MethodBasis}{《建筑防火涂料（板）工程设计、施工与验收规程》（DB11/1245-2015）}
\newcommand{\JudgeBasis}{《钢结构工程施工质量验收标准》（GB 50205-2020）}
% 检测仪器：多行用 \\ 分隔
\newcommand{\InstrumentList}{%
    GLM50+612445282 \quad 激光测距仪 \\
    MC-2000+BC200134 \quad 涂层测厚仪 \\
    （0$\sim$200）mm+87756 \quad 数显游标卡尺%
}
\newcommand{\InspStartDate}{2026年05月19日}
\newcommand{\InspEndDate}{2026年05月19日}
\newcommand{\InspEnv}{晴、32℃}
\newcommand{\InspPersonnel}{李瑾、李红月、孟林}
% 检测结论（首行自动缩进两字）
\newcommand{\Conclusion}{%
    \hspace*{2\ccwd}本次对\ProjectName{}的钢柱、钢梁进行了防火涂层厚度检测，检测结果合格。%
}
\newcommand{\Approver}{}
\newcommand{\Reviewer}{}
\newcommand{\MainInsp}{}
\newcommand{\Remarks}{/}

% —— 第1章概况正文 ——
\newcommand{\ProjectOverview}{%
\hspace*{2\ccwd}工程概况第一段描述，包括建筑基本参数。

\hspace*{2\ccwd}工程概况第二段描述，防火涂料设计要求等。

\hspace*{2\ccwd}受\ClientName{}委托，北京市建设工程质量第三检测所有限责任公司对\ProjectName{}进行了\InspItem{}。%
}

% —— 第2章检测方法正文 ——
\newcommand{\MethodDesc}{%
\hspace*{2\ccwd}按照\MethodBasis{}要求，采用测针法对防火涂层厚度进行检测。每个构件测量不少于5个点，取平均值作为该构件的厚度代表值。%
}

% ==================== 样式配置 ====================
\ctexset{
    section = {
        format     = \zihao{-3}\heiti\bfseries\centering,
        number     = \arabic{section},
        name       = {, },
        beforeskip = 0.5\baselineskip,
        afterskip  = 0.2\baselineskip
    },
    subsection = {
        format     = \zihao{4}\heiti\bfseries,
        number     = \arabic{section}.\arabic{subsection},
        name       = {, },
        beforeskip = 0.3\baselineskip,
        afterskip  = 0.1\baselineskip
    }
}
\renewcommand{\normalsize}{\zihao{-4}}

% ==================== 文档开始 ====================
\begin{document}

% ==================== 第一页：封面 ====================
\begin{titlepage}
    \pagestyle{empty}
    \heiti\sffamily\bfseries\zihao{3}

    \vspace*{15mm}

    \begin{center}
        % 字间距只用 \CJKglue，不再叠加 \hspace，避免双重间距
        {\zihao{0}\renewcommand{\CJKglue}{\hskip 0.5em}检测报告}
    \end{center}

    \vspace{5mm}

    \begin{center}
        {\kaishu\rmfamily\mdseries
        京建质检 J\textsubscript{3}—G 字 \ReportYear{}第（\ReportNo{}）号}
    \end{center}

    \vspace{55mm}

    \renewcommand{\arraystretch}{1.6}
    \begin{center}
        \kaishu\rmfamily
        % 说明：工程名称过长（>16字）时，在 \ProjectName 内部用 \\ 手动换行
        \begin{tabular}{p{3.2cm} @{\hspace{0.5em}} l}
            工\hfill 程\hfill 名\hfill 称: &
                \qquad\uline{\makebox[10.5cm][c]{\ProjectName{}}} \\
            委\hfill 托\hfill 单\hfill 位: &
                \qquad\uline{\makebox[10.5cm][c]{\ClientName{}}} \\
            检\hfill 测\hfill 类\hfill 别: &
                \qquad\uline{\makebox[10.5cm][c]{\InspType{}}} \\
        \end{tabular}
    \end{center}

    \vfill

    \begin{center}
        北京市建设工程质量第三检测所有限责任公司 \\
        \vspace{8mm}
        \ReportYear{}年\ReportMonth{}月\ReportDay{}日
    \end{center}
\end{titlepage}

% ==================== 第二页：注意事项 ====================
\clearpage
\pagestyle{empty}

\begin{center}
    \vspace*{0mm}
    {\zihao{2}\bfseries 注\quad 意\quad 事\quad 项}
    \vspace{15mm}
\end{center}

\begin{enumerate}[label=\arabic*., itemsep=1.8ex, leftmargin=2.0em]
    \zihao{-3}
    \item 本报告无"北京市建设工程质量第三检测所有限责任公司"检测专用章无效。
    \item 复制本报告无重新加盖"北京市建设工程质量第三检测所有限责任公司"检测专用章无效。
    \item 报告无主检人、审核人、批准人签字无效。
    \item 本报告涂改增删无效。
    \item 委托检验仅对来样负责。
    \item 未经本所书面批准，不得将本报告作为广告宣传使用。
    \item 若对本报告有异议，应于本报告发出之日起 15 日内向检测单位提出，逾期不受理。
\end{enumerate}

\vfill

\begin{flushleft}
    \zihao{-3}
    检测单位：北京市建设工程质量第三检测所有限责任公司 \\
    地\qquad 址：北京市西城区百万庄大街3号 \\
    邮\qquad 编：100037 \\
    电\qquad 话：010-88385491\qquad 010-88378878（管理室）\\
    传\qquad 真：010-88378878
\end{flushleft}

% ==================== 第三页：工程质量检测报告主表 ====================
\clearpage
\pagenumbering{arabic}
\pagestyle{fancy}
\fancyhf{}

% 页眉：用 \makebox 控制宽度，避免 parbox 溢出
\fancyhead[C]{%
    \zihao{-3}\mdseries\songti
    北京市建设工程质量第三检测所有限责任公司\\
    工程质量检测报告\\[0.6ex]
    \makebox[0.9\textwidth][s]{%
        京建质检 J$_3$-G 字 \ReportYear{}第（\ReportNo{}）号%
        \hfill
        第 \thepage\ 页\enspace 共 \pageref{LastPage}\ 页%
    }%
}
\cfoot{\zihao{-3}\mdseries\songti 北京市建设工程质量第三检测所有限责任公司}
\renewcommand{\headrulewidth}{1.5pt}
\renewcommand{\footrulewidth}{0pt}
\renewcommand{\footrule}{\hbox to\headwidth{\vspace{1mm}\dotfill}\vspace{0.2mm}}

\vspace*{0mm}

\noindent
{\zihao{-4}
\begin{tblr}{
  width   = \textwidth,
  colspec = {X[22,c,m] X[32,l,m] X[21,c,m] X[32,l,m] X[21,c,m] X[32,l,m]},
  hspan   = minimal,
  hlines,
  vlines,
  vline{1,Z} = {0pt},       % 去掉左右外边框，视觉上与页面融合
  row{1}     = {ht=1.5cm},
  row{2-9}   = {ht=1.0cm},
  row{10}    = {ht=1.0cm},  % 检测仪器，可多行
  row{11-12} = {ht=1.0cm},
  row{13}    = {ht=3.2cm},  % 检测结论
  row{14-15} = {ht=1.5cm},
}
工程名称 & \SetCell[c=5]{l} \ProjectName{}     & & & & \\
检测地点 & \SetCell[c=5]{l} \InspLocation{}    & & & & \\
检测项目 & \SetCell[c=5]{l} \InspItem{}        & & & & \\
施工单位 & \SetCell[c=5]{l} \ConstructUnit{}   & & & & \\
设计单位 & \SetCell[c=5]{l} \DesignUnit{}      & & & & \\
监理单位 & \SetCell[c=5]{l} \SupervUnit{}      & & & & \\
抽样依据 & \SetCell[c=5]{l} \SampleBasis{}     & & & & \\
方法依据 & \SetCell[c=5]{l} \MethodBasis{}     & & & & \\
判定依据 & \SetCell[c=5]{l} \JudgeBasis{}      & & & & \\
检测仪器 & \SetCell[c=5]{l} {\InstrumentList{}}& & & & \\
检测时间 & \SetCell[c=3]{l} \InspStartDate{}$\sim$\InspEndDate{} & & & 检测环境 & \InspEnv{} \\
检测人员 & \SetCell[c=5]{l} \InspPersonnel{}   & & & & \\
检测结论 & \SetCell[c=5]{l,t} \Conclusion{}    & & & & \\
批准人   & \Approver{}  & 审核人 & \Reviewer{} & 主检人 & \MainInsp{} \\
备注     & \SetCell[c=5]{l} \Remarks{}         & & & & \\
\end{tblr}}
 
\clearpage

% ==================== 第四页起：正文 ====================

\section{概况}

\subsection{工程概况}

\begin{spacing}{1.5}
\ProjectOverview{}
\end{spacing}


\subsection{检测及判定依据}

\begin{itemize}[label={\small\faBookOpen}, leftmargin=3em, itemsep=0.5ex]
    \item \MethodBasis{}
    \item \JudgeBasis{}
\end{itemize}


\subsection{检测人员}

\hspace*{2\ccwd}\InspPersonnel{}


\subsection{检测仪器}
\begin{center}
\vspace{0.5ex}
\noindent 表~1.4\quad 检测仪器

\vspace{0.5ex}
{\zihao{-4}
\begin{tblr}{
  width   = \textwidth,
  % 第1列为属性名，后续列为各仪器——增/减仪器时同步增/减列
  colspec = {X[2,c,m] X[3,c,m] X[3,c,m] X[3,c,m]},
  hlines, vlines,
  row{1}  = { font=\bfseries},
}
检测仪器名称   & 激光测距仪              & 涂层测厚仪               & 数显游标卡尺 \\
检测仪器编号   & GLM50+612445282         & MC-2000+BC200134         & （0$\sim$200）mm+87756 \\
检定证书编号   & JZ001号                 & JZ002号                  & JZ003号 \\
仪器使用有效期 & 2025.06$\sim$2026.06   & 2025.06$\sim$2026.06    & 2025.06$\sim$2026.06 \\
检测精度       & 1mm                     & $\pm$3\%                 & 0.01mm \\
\end{tblr}}
\end{center}


\subsection{检测环境}

\hspace*{2\ccwd}\InspEnv{}


% ==================== 第2章（根据报告类型修改章节名） ====================
\section{防火涂层厚度检测}


\subsection{检测方法}

\begin{spacing}{1.5}
\MethodDesc{}
\end{spacing}


\subsection{抽样依据}

\begin{spacing}{1.5}
\hspace*{2\ccwd}按照委托方要求，本次对\ProjectName{}的钢柱、钢梁、系杆及水平支撑进行检测。
\end{spacing}


\subsection{判定依据}

\begin{spacing}{1.5}
\hspace*{2\ccwd}依据\JudgeBasis{}及\MethodBasis{}进行判定。
\end{spacing}


\subsection{检测结果}

\noindent 表~2.4\quad 检测结果汇总

\vspace{0.5ex}
% ——————————————————————————————————————————
% 检测结果表 —— 直接在下方 tblr 内编辑数据行
%   · 列数变化：修改 colspec 中的 X[..] 数量
%   · 增/删行：复制或删除以 & 分隔的整行
%   · 判定 = max(测点) ≤ 1.5×设计厚度 且 min(测点) ≥ 0.85×设计厚度（按 GB 50205-2020）
% ——————————————————————————————————————————
{\zihao{5}
\begin{tblr}{
  width   = \textwidth,
  colspec = {X[1.2,c,m] X[1.5,c,m] X[1.8,c,m] X[1.5,c,m] X[1,c,m] X[1,c,m] X[1,c,m] X[1,c,m] X[1,c,m] X[1.2,c,m] X[1.2,c,m]},
  hlines, vlines,
  row{1-2} = {font=\bfseries},
  cell{1}{1}  = {r=2}{},
  cell{1}{2}  = {r=2}{},
  cell{1}{3}  = {r=2}{},
  cell{1}{4}  = {r=2}{},
  cell{1}{5}  = {c=5}{},
  cell{1}{10} = {r=2}{},
  cell{1}{11} = {r=2}{},
}
序号 & 构件类型 & 构件编号 & 设计厚度(mm) & 实测厚度(mm) & & & & & 平均值(mm) & 判定 \\
     &          &          &              & 1 & 2 & 3 & 4 & 5 &           &       \\
1  & 钢柱 & GZ-1 & 30 & 31 & 32 & 30 & 33 & 31 & 31.4 & 合格 \\
2  & 钢柱 & GZ-2 & 30 & 30 & 31 & 32 & 30 & 31 & 30.8 & 合格 \\
3  & 钢梁 & GL-1 & 25 & 26 & 25 & 27 & 26 & 25 & 25.8 & 合格 \\
4  & 钢梁 & GL-2 & 25 & 25 & 26 & 25 & 27 & 26 & 25.8 & 合格 \\
5  & 系杆 & XG-1 & 20 & 21 & 20 & 22 & 21 & 20 & 20.8 & 合格 \\
\end{tblr}}


% ==================== 第3章：附图 ====================
\section{附图}

% 使用方法：
%   1. 将图片文件放在 template.tex 同目录
%   2. 修改 \includegraphics 中的文件名
%   3. 复制下方图块添加更多附图

\begin{center}
    % 图片不存在时显示占位框，不影响编译
    \IfFileExists{fig_placeholder.png}{%
        \includegraphics[width=0.85\textwidth]{fig_placeholder.png}%
    }{%
        \fbox{\parbox{0.85\textwidth}{%
            \centering\vspace{2.5cm}
            【图片占位符 — 请将图片放入同目录并修改文件名】
            \vspace{2.5cm}%
        }}%
    }

    \vspace{1ex}
    图~3-1\quad 示意图（单位：cm）
\end{center}

\vfill

    \zihao{-4}（以下空白）

\end{document}